\chapter{Implementación}
\label{cap:implementacion}

Este capítulo recorre la implementación en el mismo orden que un run. Explica
cómo los módulos del mapa de la tabla~\ref{tab:mapa-sistema} convierten comandos
en hechos durables, hechos en trabajo físico y resultados físicos en evidencia
y entrega. El énfasis no está en enumerar clases, sino en las fronteras que
impiden perder estado, confundir un intento con otro o validar un commit
distinto del que finalmente recibe el usuario.

\section{Organización del código y fronteras de paquetes}

\sistema\ está implementado en TypeScript bajo el entorno de ejecución Node.js (versión 20+ ESM), estructurado como un monorepo gestionado con \texttt{pnpm workspaces}. Para asegurar la mantenibilidad y prevenir dependencias circulares, se impone una regla estricta de direccionalidad unidireccional de dependencias: \textbf{\texttt{apps $\rightarrow$ packages especializados $\rightarrow$ packages compartidos}}.

La capa \texttt{shared} contiene primitivas sin dependencias, serialización
canónica y tipos epistémicos comunes. Sobre ella se apoyan contratos y grafo;
los paquetes de dominio no dependen de Next.js, React ni de una SDK de modelos.
Así, cambiar la interfaz o el proveedor de agentes no cambia el significado de
un evento, un contrato o una evidencia. La direccionalidad se complementa con
una segunda regla: existe una sola representación canónica en cada frontera
(\texttt{SemanticPlan}, \texttt{GraphRevision}, \texttt{RunEvent} y manifiestos
Git), evitando modelos paralelos que puedan divergir.

\section{Núcleo durable: eventos, actor y efectos}
\label{sec:nucleo-durable}

\texttt{@manyhands/run-coordinator} implementa el dominio puro. Sus comandos
poseen identidad e indican la revisión esperada; sus eventos son hechos
inmutables; y \texttt{reduceRun} repliega cualquier secuencia válida en la misma
\texttt{RunProjection}. Consultar no modifica el estado y repetir un comando
idéntico devuelve el recibo previo, mientras que reutilizar su identidad con
otro contenido se rechaza.

\texttt{@manyhands/run-engine} aporta el \texttt{RunActor}. Cada run tiene un
buzón secuencial y un único escritor autorizado, aunque distintas operaciones
físicas ocurran en paralelo. El actor no ejecuta directamente Git, modelos o
procesos. Primero registra una intención durable \texttt{effect.requested}; un
adaptador realiza el efecto y produce un recibo; finalmente el actor incorpora
el resultado como evento. Esta separación evita el intervalo peligroso en el
que una operación externa ocurrió pero el sistema no dejó constancia.

\texttt{@manyhands/run-store} conserva el journal autoritativo en JSONL, con
\emph{fsync}, reemplazos atómicos, tokens de cercado y un \emph{outbox} de
efectos. Las instantáneas, índices y vistas de UI son proyecciones reconstruibles,
no una segunda verdad. Al reiniciar, el daemon reconcilia efectos no terminales
con los recibos y el estado físico antes de repetirlos. No presupone una falsa
semántica \emph{exactly once}: o demuestra que repetir es idempotente, o adopta
el resultado observado, o solicita una decisión.

\texttt{@manyhands/trace-store} guarda actividad, tiempos, consumo y diagnóstico
en un canal separado, con checksums y redacción de secretos. Una traza ayuda a
explicar qué ocurrió, pero nunca puede promover un intento ni cambiar el ciclo
de vida: esa autoridad pertenece exclusivamente a los eventos del run.

\section{Disponibilidad, priorización y conducción de olas}

\texttt{@manyhands/scheduler} evalúa continuamente la disponibilidad de cada
nodo. Comprueba contratos aprobados, artefactos frescos, decisiones resueltas,
autoridad de recursos, capacidad del ejecutor, presupuesto y concesiones
físicas. El resultado no es una fase global, sino una frontera de nodos
\emph{hard-ready} que cambia después de cada evento.

\texttt{@manyhands/conflict-risk} estima interferencias por rutas, símbolos y
dependencias para priorizar dentro de esa frontera. Su salida es consultiva:
puede reducir concurrencia o cambiar el orden, pero no inventar dependencias ni
autorizar escrituras. \texttt{@manyhands/orchestrator-graph} conduce el bucle:
selecciona una ola compatible, solicita intentos, incorpora artefactos adoptados
y vuelve a calcular la frontera. De ese modo hojas e integraciones intermedias
pueden alternarse sin esperar a que termine todo el grafo.

\section{Intentos, artefactos y ejecución aislada}
\label{sec:workers-aislados}

Cada ejecución, reparación o integración crea un intento inmutable. Su
\texttt{InputFingerprint} resume nodo, contratos, árbol base, vista del
repositorio, manifiestos consumidos, ejecutor y capacidad de sandbox. Si cambia
una entrada, se crea otro intento y la evidencia anterior queda obsoleta; un
intento terminal nunca se reabre.

\texttt{ExecutionBaseBuilder} materializa un árbol exacto desde el commit base
y únicamente los manifiestos declarados. El agente edita, pero no decide la
identidad del resultado: el orquestador observa el diff, verifica alcance y
autoridad, y recién entonces crea el commit candidato. \texttt{ArtifactBuilder}
extrae objetos Git y los describe en manifiestos inmutables que conservan
archivos binarios, permisos, enlaces y eliminaciones. El estado de adopción o
validación se registra aparte, para no mutar la identidad del artefacto.

La ejecución de código generado plantea riesgos de corrupción del sistema de
archivos y desbordes de recursos. Por eso se aplican capas independientes de
aislamiento y custodia:

\subsection{Pool de Git worktrees y reciclado determinista}

Cada intento asignado a un agente se despacha en un worktree independiente dentro de un pool administrado por el orquestador. Antes de iniciar la ejecución, el worktree se restablece al commit base requerido mediante una secuencia atómica de comandos Git:

Cada slot se materializa desde una base Git identificada y se entrega a un solo
intento bajo una concesión exclusiva. Cuando el agente finaliza, el orquestador
inspecciona el árbol mediante \texttt{git diff} y
\texttt{git status --porcelain=v1}. La salida textual del modelo nunca se usa
para decidir qué archivos cambiaron: \textbf{la fuente física de verdad es el
diff del worktree}. Antes de reutilizar un slot, el gestor comprueba que no
existan operaciones Git incompletas ni propietarios todavía activos.

\subsection{Verificación estricta de alcance en tiempo de ejecución (\emph{Scope Guard})}

Para evitar que un agente modifique archivos ajenos a su sub-tarea, el subsistema de ejecución aplica un validador de alcance en dos fases con política de \emph{denegación prioritaria} (\emph{deny-first}):

\begin{enumerate}
  \item \textbf{Normalización POSIX y prevención de escapes:} Cada ruta tocada se normaliza resolviendo enlaces simbólicos (\emph{symlinks}) y secuencias de ascenso de directorios (\texttt{..}). Cualquier intento de escapar de la raíz del worktree genera un rechazo inmediato.
  \item \textbf{Verificación contra el contrato de alcance:} Todo archivo modificado, creado o eliminado en el diff debe pertenecer al conjunto de rutas permitidas (\texttt{writePaths}) del \texttt{ScopeContract}. Si el agente toca un solo archivo no autorizado (por ejemplo, \texttt{package.json} o archivos de configuración global cuando solo tenía permiso sobre \texttt{src/domain/}), el intento se aborta de inmediato clasificándose con el estado terminal \texttt{scope\_violation}, y ningún commit es adoptado por el orquestador.
\end{enumerate}

\subsection{Supervisión de procesos, comunicación IPC y saneamiento de entorno}

\begin{itemize}
  \item \textbf{Supervisión del árbol de procesos del SO:} La CLI del agente se inicia como un proceso supervisado. En Windows, un \emph{Job Object} nativo mantiene la custodia del árbol completo; si vence un plazo o el usuario cancela, el cierre del trabajo termina también los descendientes. El efecto sólo se considera físicamente final después de comprobar la identidad y el resultado del proceso supervisado.
  \item \textbf{Sandbox y credenciales acotadas:} El sandbox de ejecución limita la escritura al worktree. La credencial necesaria para invocar al modelo se copia a un hogar efímero identificado por run e intento y se revoca después del resultado físico terminal. Así, el worker no necesita heredar el entorno completo del daemon ni conservar autenticación entre intentos.
  \item \textbf{Concesiones con testigo (\emph{Witness Leases}):} Toda mutación sobre un repositorio o estado del run exige una concesión durable con vigencia demostrable \cite{burrows2006chubby}. Si un worker tarda demasiado o el daemon se reinicia, cualquier resultado tardío se descarta.
\end{itemize}

Un worktree separa checkouts; no constituye por sí solo un sandbox. La
restricción de archivos, la custodia de procesos, la política de red y el
alcance de credenciales son capacidades distintas y se verifican por separado.
En Windows, \texttt{windows-job-runner} aporta la custodia nativa de procesos y
\texttt{windows-ipc-acl} protege los \emph{Named Pipes} mediante DACLs acotadas
al usuario y al sistema, incluyendo defensas frente a redirecciones del sistema
de archivos.

\section{Integración ascendente y reparación semántica de conflictos}
\label{sec:integracion-ascendente}

La integración es un intento de primer orden modelado en el grafo. Los nodos
compuestos no son agrupadores visuales: construyen un candidato combinado desde
los manifiestos frescos de sus hijos y, cuando el contrato lo prevé, realizan
ediciones sólo sobre recursos compartidos propiedad del padre:

\begin{enumerate}
  \item \textbf{Composición en árbol limpio:} El nodo compuesto materializa una base exacta e integra los cambios declarados de sus hijos siguiendo el orden topológico de artefactos, verificando preimágenes y versiones.
  \item \textbf{Verificación de costuras e interfaces:} Una vez fusionados los cambios, el compuesto ejecuta las obligaciones de validación de nivel integración (\texttt{layer: "integration"}) asociadas a sus \texttt{SeamContract}s. Que dos ramas apliquen sin conflictos textuales de Git no garantiza corrección semántica: la integración solo se considera exitosa si las pruebas de integración conjuntas pasan.
  \item \textbf{Reparación semántica de conflictos (*Semantic Conflict Repair*):} Si la composición o una prueba falla, el orquestador crea un nuevo intento con los diagnósticos y el contexto del padre. Una corrección necesaria sobre recursos de un hijo vuelve a ese hijo o exige una enmienda; la integración no recibe permisos universales.
\end{enumerate}

\section{Recuperación causal y decisiones durante el run}

Los fallos no comparten un contador universal de reintentos. El dominio los
clasifica por causa: entrada inválida, entorno o proveedor transitorio, violación
de alcance, proceso perdido, incompatibilidad de integración o evidencia
insuficiente, entre otras. Reintentar es válido sólo si cambia una condición
relevante, por ejemplo el contexto, la capacidad disponible o el diagnóstico
entregado al agente. Así se evita repetir indefinidamente la misma ejecución.

La recuperación conserva el trabajo ya adoptado cuyos fingerprints continúan
vigentes. Un error de una hoja no obliga a reiniciar el run: puede originar una
reparación local; una modificación de contrato genera una revisión o
\texttt{Amendment} inmutable y vuelve obsoletos únicamente los descendientes
afectados. Cuando la resolución requiere criterio humano, una \texttt{Decision}
declara los nodos alcanzados y permite que las ramas independientes sigan
avanzando.

\section{Matriz de Evidencias (Evidence Matrix)}
\label{sec:matriz-evidencias}

La validez de un resultado de software en ManyHands se sintetiza en la \textbf{Matriz de Evidencias} (\texttt{EvidenceMatrix}). Cada celda de la matriz vincula formalmente:

\begin{multline}
\text{CeldaEvidencia} = \big(\text{CriterioId},\;
\text{ObligacionValidacionId},\; \text{CommitSha},\\
\text{Resultado},\; \text{TrazaComando}\big)
\label{eq:celda-evidencia}
\end{multline}

\noindent La matriz aplica las siguientes reglas de auditoría:

\begin{itemize}
  \item \textbf{Prohibición de pruebas debilitadas (*Test Integrity Guard*):} El verificador de evidencias audita los diffs de pruebas generados por los agentes. Si detecta la eliminación de aserciones existentes, la adición de modificadores para omitir pruebas (como \texttt{it.skip}, \texttt{describe.only} o \texttt{xit}) o la modificación de configuraciones del runner de pruebas para ignorar fallos, el resultado de validación es clasificado como \texttt{fraudulent\_evidence} y rechazado de inmediato.
  \item \textbf{Controles negativos cuando son aplicables:} Una prueba nueva puede ejecutarse sobre el commit base para comprobar que detecta la ausencia del cambio. Si pasa también sobre el código previo, no sirve como evidencia discriminante para esa obligación.
  \item \textbf{Política contra pruebas inestables (*Flaky Test Policy*):} Las obligaciones de validación críticas declaran \texttt{flakyPolicy: "forbid"}. Si una prueba exhibe intermitencia no determinista durante repeticiones consecutivas sobre el mismo commit, la matriz clasifica la celda como inestable y bloquea la promoción del artefacto.
\end{itemize}

Un candidato sólo alcanza el estado \texttt{verified} cuando todas las
obligaciones requeridas están satisfechas sobre ese candidato. Las observaciones
advisory o inconclusas pueden conservarse para diagnóstico, pero no sustituyen
una obligación bloqueante.

\section{Publicación y Entrega Transaccional (TransactionalDeliveryPublisher)}
\label{sec:entrega-transaccional}

La fase final de entrega es gestionada por la clase \texttt{TransactionalDeliveryPublisher} en \texttt{@manyhands/execution-core}. La publicación garantiza atomicidad y trazabilidad mediante el protocolo de tres fases ilustrado en la figura~\ref{fig:entrega-transaccional}:

\begin{figure}[htbp]
  \centering
  \begin{tikzpicture}[
    box/.style={draw=black!70, fill=black!3, rounded corners=3pt, align=center, font=\small, minimum height=1.0cm, text width=3.4cm},
    act/.style={draw=blue!80!black, fill=blue!5, thick, rounded corners=3pt, align=center, font=\small, minimum height=1.0cm, text width=3.4cm},
    term/.style={draw=green!60!black, fill=green!8, thick, rounded corners=3pt, align=center, font=\small, minimum height=1.0cm, text width=3.4cm},
    arrow/.style={-{Stealth[length=2.5mm]}, thick, draw=black!70}
  ]
    \node[box] (prep) {\textbf{1. Preparación}\\ Construcción del Manifiesto Final (\texttt{FinalManifest})};
    \node[act, right=1.0cm of prep] (val) {\textbf{2. Validación exacta}\\ Comparación de SHA, árbol Git y Evidence Matrix};
    \node[term, right=1.0cm of val] (pub) {\textbf{3. Publicación}\\ Fast-forward atómico y emisión de recibo};

    \draw[arrow] (prep) -- (val);
    \draw[arrow] (val) -- (pub);
  \end{tikzpicture}
  \caption{Protocolo de publicación transaccional de entrega.}
  \label{fig:entrega-transaccional}
\end{figure}

\begin{enumerate}
  \item \textbf{Manifiesto de artefacto final (\texttt{FinalArtifactManifest}):} Documento inmutable que contiene el SHA del commit candidato, el SHA del árbol físico de Git (\texttt{treeSha}), el ID de la Matriz de Evidencias verificada y la rama de destino.
  \item \textbf{Validación de concurrencia y precondiciones:} El publicador inspecciona el repositorio de destino y comprueba que la rama destino no haya avanzado concurrentemente (\texttt{targetHead == targetHeadBefore}) y que el directorio de trabajo esté limpio.
  \item \textbf{Publicación y Recibo de Entrega (\texttt{TransactionalDeliveryReceipt}):} El publicador ejecuta un avance rápido atómico (\texttt{git merge --ff-only <finalSha>}) y emite un recibo inmutable persistido en el journal de eventos, conteniendo las identidades criptográficas de antes y después del cambio.
\end{enumerate}

\section{Persistencia append-only basada en eventos y Daemon local}

El estado se persiste mediante el \emph{Event Store} de
\texttt{@manyhands/run-store} en formato JSON Lines, con semántica de sólo
anexión. Cada lote se publica de manera segura:

\begin{equation}
\text{archivo\_temp} \xrightarrow{\text{flush a disco}} \text{fsync()} \xrightarrow{\text{reemplazo atómico}} \text{archivo\_definitivo}
\label{eq:atomic-write}
\end{equation}

El \textbf{daemon local} es el anfitrión privilegiado y único dueño de actores,
journals, temporizadores, credenciales intermediadas y workers. Una concesión de
instalación impide que dos daemons posean el mismo directorio de estado. En el
arranque, reconcilia journals y efectos pendientes antes de aceptar comandos.

\texttt{apps/web} no se conecta desde JavaScript al canal privilegiado. Su
servidor Next.js funciona como BFF: valida \texttt{Host}, \texttt{Origin}, token
de sesión y tipo de contenido, y recién entonces traduce comandos y consultas a
IPC local autenticado con HMAC, nonce y protección contra repetición. En
Windows, el transporte usa un \emph{Named Pipe} con ACL protegida. Si el
navegador se cierra o recarga, el daemon sigue ejecutando y la vista se
reconstruye desde la proyección durable.

\section{Interfaz web de supervisión y accesibilidad}

La interfaz gráfica está construida sobre Next.js y React Flow como un
\textbf{Command Center}. Los comandos cambian el mundo; las consultas y el
stream ordenado de eventos sólo lo observan. SSE conserva el número de secuencia
y, tras una desconexión o un hueco, vuelve a consultar la proyección antes de
continuar. Por eso una recarga no crea otra verdad ni depende de memoria local:

\begin{itemize}
  \item \textbf{Estabilidad visual del lienzo:} Se implementa el invariante estricto de que \textbf{el lienzo del grafo nunca se recentra, reencuadra ni altera su nivel de zoom de forma automática} ante la llegada de eventos de ejecución o cambios de estado de los nodos. El desarrollador mantiene el control absoluto de su encuadre de trabajo.
  \item \textbf{Cola global de decisiones no bloqueante:} Las decisiones humanas pendientes se presentan en una cola accesible. El trabajo en ramas independientes del grafo continúa ejecutándose concurrentemente: una decisión pendiente detiene únicamente aquellos nodos cuya disponibilidad depende directamente del resultado de dicha decisión.
  \item \textbf{Accesibilidad y movimiento reducido:} La interfaz fue diseñada siguiendo WCAG 2.2 AA \cite{w3c2023wcag22}: utiliza nombres accesibles, foco visible, contraste suficiente y respeta \texttt{prefers-reduced-motion}. Estas decisiones permiten supervisar el sistema sin depender exclusivamente del color o de animaciones.
  \item \textbf{Estados honestos y observabilidad:} La UI distingue candidato, verificado, obsoleto, fallido y entregado, y vincula cada nodo con su actividad y diagnóstico. Costos y trazas ayudan a operar el run, pero no se presentan como prueba de corrección.
\end{itemize}
